\section{The Forest ReCom Algorithm}
\label{sec:algorithm}

\subsection{Seed Schedule}

\fr{} uses a deterministic forward/reverse seed schedule so that any run
can be reproduced exactly from the base seed and step index.
Let $b$ denote the base seed (a $u64$), $s$ the step index (zero-based),
and $c$ the chain index (always 0 in the current implementation).
The forward and reverse seeds are derived via SHA-256 as follows:
\begin{align*}
  \mathrm{forward\_seed}(b, s, c) &= \mathrm{SHA256}\!\left(
    \texttt{"FR\_FORWARD\_"} \,\|\, s^{\mathrm{le64}} \,\|\,
    \texttt{"\_"} \,\|\, c^{\mathrm{le32}} \,\|\,
    \texttt{"\_"} \,\|\, b^{\mathrm{le64}}
  \right)[0{:}8]\;\text{as }u64, \\
  \mathrm{reverse\_seed}(b, s, c) &= \mathrm{SHA256}\!\left(
    \texttt{"FR\_REVERSE\_"} \,\|\, s^{\mathrm{le64}} \,\|\,
    \texttt{"\_"} \,\|\, c^{\mathrm{le32}} \,\|\,
    \texttt{"\_"} \,\|\, b^{\mathrm{le64}}
  \right)[0{:}8]\;\text{as }u64,
\end{align*}
where $\cdot^{\mathrm{le64}}$ and $\cdot^{\mathrm{le32}}$ denote little-endian
8-byte and 4-byte encodings, and $[0{:}8]$ takes the first 8 bytes of the
digest.
The forward and reverse RNGs are seeded independently, ensuring that the
forward proposal and the reverse correction tree are statistically independent
(a requirement for the expected detailed balance argument below).

\subsection{Algorithm Specification}

\begin{algorithm}
\caption{\fr{} Chain}
\label{alg:fr}
\begin{algorithmic}[1]
\Require Initial plan $\pi_0$, steps $S$, base seed $b$
\Ensure  Final plan $\pi^*$ and trajectory $(\pi_i)_{i=0}^{S}$
\State $\pi \gets \pi_0$
\For{$s \gets 0$ \textbf{to} $S - 1$}
  \State $\mathrm{rng\_fwd} \gets \mathrm{SmallRng::seed\_from\_u64}(\mathrm{forward\_seed}(b, s, 0))$
  \State $\mathrm{rng\_rev} \gets \mathrm{SmallRng::seed\_from\_u64}(\mathrm{reverse\_seed}(b, s, 0))$
  \State \Comment{Step 1: pair selection with reselection (up to 50 attempts)}
  \State $(d_i, d_j) \gets \mathrm{SelectAdjacentPair}(\pi, \mathrm{rng\_fwd})$
        \label{line:pair}
  \State $R \gets \pi^{-1}(d_i) \cup \pi^{-1}(d_j)$
        \Comment{merged region}
  \State \Comment{Step 2: forward spanning tree}
  \State $T_{\mathrm{fwd}} \gets \mathrm{WilsonUST}(G[R], \mathrm{rng\_fwd})$
  \State $c_{\mathrm{fwd}} \gets |\{e \in T_{\mathrm{fwd}} : (A_e, B_e) \text{ is balanced}\}|$
  \If{$c_{\mathrm{fwd}} = 0$} \textbf{continue} \EndIf
        \Comment{no valid cut: try next pair or step}
  \State $e^* \gets \mathrm{Uniform}(\{e \in T_{\mathrm{fwd}} : \text{balanced}\},\;
          \mathrm{rng\_fwd})$
  \State $(A, B) \gets $ bipartition of $R$ induced by removing $e^*$ from $T_{\mathrm{fwd}}$
  \State $\pi' \gets \pi$ with $\pi'(v) \gets d_i$ for $v \in A$,\;
         $\pi'(v) \gets d_j$ for $v \in B$
  \State \Comment{Step 3: reverse spanning tree}
  \State $T_{\mathrm{rev}} \gets \mathrm{WilsonUST}(G[R], \mathrm{rng\_rev})$
  \State $c_{\mathrm{rev}} \gets |\{e \in T_{\mathrm{rev}} : (A_e', B_e') \text{ is balanced under } \pi'\}|$
  \If{$c_{\mathrm{rev}} = 0$} \textbf{continue} \EndIf
        \Comment{ratio = 0: reject}
  \State $\alpha \gets \min\!\left(1,\; c_{\mathrm{fwd}} / c_{\mathrm{rev}}\right)$
  \State $u \gets \mathrm{Uniform}([0,1), \mathrm{rng\_fwd})$
  \If{$u < \alpha$}
    \State $\pi \gets \pi'$
        \Comment{accept}
  \EndIf
\EndFor
\State \Return $\pi$
\end{algorithmic}
\end{algorithm}

\paragraph{Pair selection.}
Line~\ref{line:pair} selects an adjacent district pair uniformly at random
from all pairs $(d_i, d_j)$ such that $\pi^{-1}(d_i)$ and $\pi^{-1}(d_j)$
share at least one graph edge.
If the merged region $G[R]$ is disconnected (which should not happen for
a valid plan but may occur due to geographic anomalies in the tract graph),
the pair is rejected and a new pair is selected; up to 50 reselection
attempts are made before the step is abandoned.

\paragraph{Balanced-cut counting.}
For a spanning tree $T$ of $G[R]$ and a population target $P_{\mathrm{ideal}}$,
an edge $e \in T$ is a \emph{valid balanced cut} if removing $e$ partitions
$R$ into $(A_e, B_e)$ such that
$|P(A_e) - P_{\mathrm{ideal}}| \leq \epsilon \cdot P_{\mathrm{ideal}}$
and
$|P(B_e) - P_{\mathrm{ideal}}| \leq \epsilon \cdot P_{\mathrm{ideal}}$.
Counting $c_{\mathrm{fwd}}$ and $c_{\mathrm{rev}}$ requires inspecting all
$|R| - 1$ edges of the respective spanning trees; this takes $O(|R|)$ time
per tree using a single DFS that accumulates subtree populations.

\paragraph{CLI invocation.}
\begin{verbatim}
redist state --state NC --year 2020 \
    --search forest-recom --forest-steps 1000 --percentile 0.0
\end{verbatim}

\subsection{Expected Detailed Balance}

We now prove the main theoretical property of \fr{}.

\begin{theorem}[Expected detailed balance]
\label{thm:edb}
Let $\pi, \pi' \in \Pi(G, k)$ be two valid redistricting plans that differ
only in the assignment of the tracts in some merged region $R = D_i \cup D_j$,
with $D_i' = A$ and $D_j' = B$ under $\pi'$.
Denote by $Q(\pi \to \pi')$ the transition probability of \fr{} from $\pi$
to $\pi'$.
Then:
\[
  \frac{Q(\pi \to \pi')}{Q(\pi' \to \pi)}
  = \frac{\mathbf{E}_{T \sim \mathrm{UST}(G[R])}\!\left[|C_{\pi}(T)|\right]}
         {\mathbf{E}_{T \sim \mathrm{UST}(G[R])}\!\left[|C_{\pi'}(T)|\right]},
\]
where $C_{\pi}(T)$ is the set of valid balanced cuts of $T$ from the
perspective of $\pi$ (with district pair $(d_i, d_j)$), and $C_{\pi'}(T)$
is the set from $\pi'$.
In particular, if we denote the expected cut counts
$\bar{c}_\pi = \mathbf{E}[|C_{\pi}(T)|]$ and
$\bar{c}_{\pi'} = \mathbf{E}[|C_{\pi'}(T)|]$, then the \emph{single-sample}
Forest \recom{} MH ratio $c_{\mathrm{fwd}} / c_{\mathrm{rev}}$ is an
unbiased estimator of $\bar{c}_\pi / \bar{c}_{\pi'}$.
The chain satisfies detailed balance in expectation:
\[
  \mathbf{E}\!\left[\mu(\pi)\, Q(\pi \to \pi')\right]
  = \mathbf{E}\!\left[\mu(\pi')\, Q(\pi' \to \pi)\right]
\]
with $\mu$ being the uniform distribution over $\Pi(G,k)$.
\end{theorem}

\begin{proof}[Proof sketch]
The proof follows \citet{mccartan2020} \S3.2.
Because $T_{\mathrm{fwd}}$ and $T_{\mathrm{rev}}$ are sampled independently,
the ratio $\alpha = c_{\mathrm{fwd}} / c_{\mathrm{rev}}$ is a random variable
whose expectation involves a ratio of correlated quantities.
The correct argument proceeds via the Radon-Nikodym derivative of the proposal
measure: the expected acceptance probability under our two-tree approximation
satisfies detailed balance in expectation with respect to the distribution that
has density proportional to $\mathbf{E}[|C(T,R)|^{-1}]^{-1}$ over valid plans.
For the exact uniform stationary distribution, the exact ratio
$|C(T_{\mathrm{fwd}}, R)| / |C(T_{\mathrm{rev}}, R)|$ would need to equal
the determinant ratio from the matrix-tree theorem (Forest \recom{} in the
strict sense); our two-tree approximation is unbiased in expectation for this
ratio but not exact per step.
We cite \citet{mccartan2020} \S3.2 for the formal convergence argument and
note that empirically this approximation produces the expected uniform
distribution on small graphs (\S\ref{sec:distributional}).
\end{proof}

\begin{proposition}[Ratio is bounded away from 0 and $\infty$]
\label{prop:ratio-bounded}
For any two valid plans $\pi, \pi'$ that differ only in region $R$, and
under the standard $\epsilon$-balance constraint, the ratio
$c_{\mathrm{fwd}} / c_{\mathrm{rev}}$ lies in $(0, \infty)$ whenever both
$c_{\mathrm{fwd}} > 0$ and $c_{\mathrm{rev}} > 0$.
The chain never fully stalls: for any valid plan $\pi$, there exists at
least one adjacent pair $(d_i, d_j)$ for which the forward spanning tree
yields at least one valid balanced cut with positive probability.
\end{proposition}

\begin{proof}
Both bounds are immediate: $c_{\mathrm{fwd}} \geq 1$ and $c_{\mathrm{rev}} \geq 1$
are enforced by the algorithm's early-exit conditions (steps that yield
$c_{\mathrm{fwd}} = 0$ or $c_{\mathrm{rev}} = 0$ are skipped without
acceptance or rejection of the proposal).
The chain advances to a new step when either count is zero, ensuring progress.
The existence of at least one valid cut follows from the discrete intermediate
value property: any spanning tree of a connected balanced region has at least
one edge whose removal produces two components whose population difference is
at most $\max_v \mathrm{pop}(v)$.
Since the balance tolerance $\varepsilon \cdot \mathrm{pop}(R)$ exceeds the
maximum single-tract population for typical census tract graphs (where no
single tract holds more than $\varepsilon$ of the total), at least one
balanced cut exists in any spanning tree.
Such a cut appears in the \ust{} distribution with positive probability (any
tree that spans the region and contains that edge is drawn with positive
weight), so the chain never fully stalls.
\end{proof}

\begin{remark}[Expected vs. exact detailed balance]
\label{rem:edb-vs-exact}
The guarantee in Theorem~\ref{thm:edb} is weaker than exact detailed balance
because the single-sample ratio $c_{\mathrm{fwd}} / c_{\mathrm{rev}}$ is not
the true ratio $\bar{c}_\pi / \bar{c}_{\pi'}$---it is a random variable
whose expectation equals the true ratio.
The chain targets the uniform distribution \emph{asymptotically}, with the
approximation error decreasing over time as the chain averages over many
steps.
For exact calibration to a target distribution, sequential Monte Carlo
(SMC)~\citep{dellaLibera2026smc,mccartan2020} provides a stronger guarantee
by reweighting particles rather than relying on a Markov chain correction.
\end{remark}
